\section{Conclusion}

In this paper, we present ProtoDRIFT, a prototype-disentangled, deviation-referenced invariant forecasting framework for cross-city anomalous urban flow prediction. Unlike existing methods that model anomalies directly in the absolute-value space, ProtoDRIFT decomposes forecasting into city-specific normal reference retrieval and deviation-space modeling, thereby stripping transferable anomaly semantics away from city-specific scales. By learning a city-agnostic anomaly prototype bank in the deviation space, ProtoDRIFT distills reusable anomalous patterns into compact and semantically interpretable representations, and explicitly separates transferable anomaly semantics from city-specific responses through mutual information constraints and a city-conditioned branch. In bidirectional zero-shot cross-city transfer experiments between NYC and CHI, ProtoDRIFT achieves leading performance on both the full and anomaly subsets, and the ablation and visualization analyses further verify the contribution of each component as well as the city-agnostic nature and interpretability of the anomaly prototypes, demonstrating that deviation-space prototype disentanglement is an effective route to cross-city transfer of anomaly knowledge. In future work, we will explore stronger invariance constraints under multi-city joint training, introduce weakly supervised anomaly labels to reduce the reliance on anomaly masks, and extend the framework to more general urban spatio-temporal forecasting tasks such as travel demand and road speed prediction.